% ===================================================================
%  HOTO Na 11 — Gari da gine-ginensa. --- Gobara.
%
%  Traduction, faite en 2026, en haoussa d'aujourd'hui, sur l'ido de
%  texto/io. Regles de la colonne : voir 10-hoto-01.tex. Deux scenes,
%  deux numerotations. Ouverture de la TROISIEME SERIE : « JERI NA
%  UKU ». Deux blocs marques « ¶2 » par ordo.py.
%
%  C'EST LE TABLEAU QUI N'AVAIT AUCUN GROS PLAN DANS LES QUARANTE-DEUX
%  COLONNES, ET IL EN A MAINTENANT 89. La planche numerote c1/c2 pour
%  ses deux scenes mais grave ses numeros NUS ; html.py cherchait
%  « c1:7 » et ne trouvait rien, dans toutes les langues a la fois. On
%  ne s'en est apercu qu'en comptant les boutons tableau par tableau,
%  pendant l'ecriture de la colonne ourdoue.
%
%  LA VILLE HAOUSSA A DEUX AGES, COMME LA PERSANE ET LA JAVANAISE AU
%  MEME TABLEAU. Ce qui existait avant le XIXe siecle porte un nom
%  d'ici : birni la ville murée, gari la ville, unguwa le quartier,
%  ganuwa l'enceinte, ƙofar gari la porte (8), dandali la place (96),
%  hasumiya la tour (9), kasuwa le marche, masallaci, makabarta,
%  shirayi le peristyle (78). Ce que l'Etat moderne a bati porte un
%  nom anglais : banki (32), gidan waya (43), tiram (31), asibiti
%  (39), bariki (16), kwastan (10).
%
%  MAIS L'EXCEPTION DE LA DOUANE NE SE PRODUIT PAS ICI, ET C'EST LA
%  PREMIERE FOIS. Au tableau 11 javanais, « pabean » etait le seul
%  batiment moderne a garder un nom d'ici ; le persan faisait la meme
%  exception avec « گمرک », mille ans dans la langue. Deux colonnes,
%  deux mots differents, une seule raison : la douane est le plus
%  vieux des batiments modernes, parce qu'on percoit au port depuis
%  toujours. Le haoussa ne suit pas. Il n'a jamais eu de douane de
%  PORT : le commerce transsaharien se taxait a la porte de la ville,
%  et cela s'appelait « fito » ou « kuɗin ƙofa », le prix de la porte.
%  Le batiment de la planche, sur un quai, entre un fleuve et une
%  ruelle, est arrive avec les Britanniques et avec son nom. UNE
%  EXCEPTION QUI TIENT DEUX FOIS ET TOMBE LA TROISIEME EN DIT PLUS
%  QU'UNE QUI TIENDRAIT TOUJOURS : elle ne tenait pas a la douane,
%  elle tenait au PORT.
%
%  ET « ganuwa » RENDS « rempari » SANS UN MOT DE TROP. L'ido dit
%  qu'on trouve encore quelques restes des remparts de la ville, avec
%  la porte de l'ancien chateau fort flanquee de ses deux tours. La
%  ganuwa de Kano fait quatorze kilometres de terre battue, elle a
%  quinze portes, elle est debout, et l'on entre par la ƙofar gari
%  comme on y entrait au XVe siecle. Aucun mot a chercher, aucune
%  chose a transposer : c'est le meme objet.
%
%  LE VIOLON (81) NE PEUT PAS S'ECRIRE, ET C'EST L'ALPHABET QUI
%  TRANCHE. Le boko n'a pas de v ; le mot doit donc se refaire, et il
%  se refait en b, comme vitamin donne bitamin : « bayolin ». Le
%  haoussa a par ailleurs « goge », la vielle a une corde que frotte
%  le maroƙi, et goge reste ou il est. C'est exactement la solution
%  persane — ویولن pour l'instrument d'orchestre, کمانچه pour le
%  sien — mais elle n'a pas ete choisie : elle a ete IMPOSEE par une
%  lettre qui manque. Le releve compte maintenant cinq reponses au
%  violon : سارنگی refuse en ourdou, biola pris au portugais dans les
%  deux colonnes de Java, ویولن accepte a cote de کمانچه, bayolin
%  refait au b. Le tambour (85), lui, est « ganga » sans un instant
%  d'hesitation : un tambour est un tambour.
%
%  LA POSTE (43) EST « gidan waya », LA MAISON DU FIL. C'est le
%  telegraphe qui a nomme la poste, et non l'inverse — le fil est
%  arrive au Nigeria avant les timbres. Le facteur (45) est « mai kai
%  wasiƙa », celui qui porte la lettre, avec un mot arabe pour la
%  lettre et un mot haoussa pour le porter.
%
%  ET « keke » ATTEINT SEPT : les anciennes voitures de tramway
%  tirees par des chevaux (31) sont des « kekunan doki ». Bicyclette,
%  voiture attelee, rouet, roue de moulin, funiculaire, roue a aubes,
%  tramway a chevaux.
%
%  ENFIN L'INCENDIE RETOURNE AU HAOUSSA D'UN SEUL COUP, comme il
%  retournait au javanais et au persan aux memes alineas : wuta le
%  feu, gobara l'incendie (74), hayaƙi la fumee (75), harshen wuta la
%  flamme, tsani l'echelle (72), kashe gobara eteindre. Une ville
%  brule dans la langue de ceux qui l'eteignent, meme quand ses
%  batiments portent des noms etrangers. Quatrieme fois que ce
%  partage se verifie.
% ===================================================================

%%K t11-apar-1 apar
\VUsaut{2.0mm}
\VUcentre{13.2pt}{90}{JERI NA UKU}
\VUsaut{3.0mm}
\VUfilet{16mm}
\VUsaut{5.6mm}
\VUtitre{15.0pt}{60}{HOTO Na 11}
\VUsaut{1.4mm}
\VUpk{11.4pt}{40}{Gari da gine-ginensa. --- Gobara.}
\VUsaut{2.2mm}

%%K t11-tit-1 sub
\VUpk{10.2pt}{40}{Unguwar bayan gari.}
\VUsaut{3.0mm}

%%K t11-c1-01-1 p
1. --- Ina zaune a wani babban gari da yake kilomita 100 daga Babban
Teku, wanda duk da haka yake \VUgras{tashar jiragen ruwa}
\textsuperscript{(1)} ta matsayi mai girma. Kogin da yake iyakance shi
yana da faɗin mita 400 kuma yana yin kyakkyawan mashigi.

%%K t11-c1-02-1 p
2. --- Dukan ɓangaren garin da yake kan \VUgras{tashoshin ruwa}
\textsuperscript{(2)} yana da kamannin teku da masana’antu gaba ɗaya, yana
kuma yin \VUgras{unguwar bayan gari} \textsuperscript{(3)} wadda matuƙan
jirgin ruwa da ma’aikata kaɗai suke zaune a ciki. Waɗannan suna aiki a
\VUgras{masana’antu} \textsuperscript{(4)} da \VUgras{gidan iskar gas}
\textsuperscript{(5)} wanda ake ganin doguwar \VUgras{bututun hayaƙinsa}
\textsuperscript{(6)} da \VUgras{ma’aunin gas} daga nesa ƙwarai.

%%K t11-c1-03-1 p
3. --- A tsakiyar zamanai, an yi wa garin ganuwa, har yanzu kuwa akan
sami waɗansu saurar ganuwarsa, musamman \VUgras{ƙofar gari}
\textsuperscript{(8)} ta tsohuwar \VUgras{kagara} \textsuperscript{(7)} wadda
\VUgras{hasumiyoyinta} \textsuperscript{(9)} biyu suke gefenta, waɗanda yanzu
ake amfani da su a matsayin \VUgras{kurkuku.}

%%K t11-c1-04-1 p
4. --- A dukan wannan \VUgras{unguwa,} wadda take da kamannin tsufa
ƙwarai, gini ɗaya kaɗai yake na zamani : shi ne \VUgras{gidan kwastan}
\textsuperscript{(10)}. Yana tsakanin tashar ruwa da \VUgras{ƙaramar hanya}
\textsuperscript{(11)} wadda take kai wa tsohuwar \VUgras{majalisar gari}
\textsuperscript{(12)}. Ana gane wannan gini na ƙarshe cikin sauƙi saboda
babbar \VUgras{hasumiyar ƙararrawa} \textsuperscript{(13)} wadda take hangen
tsohon garin.

%%K t11-c1-05-1 p
5. --- A ɗaya gefen hanyar, akwai gini \VUgras{a tsaye} da aka gina da
salo ɗaya da gidan kwastan : shi ne \VUgras{Kasuwar hannun jari}
\textsuperscript{(14)}. Ɗaya daga cikin fuskokinsa yana fuskantar ƙaramin
fili inda \VUgras{gidan gwamna} \textsuperscript{(15)} yake. Hakika, garinmu,
ba tare da ya zama babban birni ba, duk da haka babban birnin lardi ne
mai muhimmanci.

%%K t11-tit-2 sub
\VUsaut{5.0mm}
\VUpk{10.2pt}{40}{Ɓangaren Gari da Ya Fi Tsawo.}
\VUsaut{3.4mm}

%%K t11-c1-06-1 p
6. --- Sansanin sojoji, mai yawa sosai, yana zaune a manyan
\VUgras{barikin sojoji} \textsuperscript{(16)} masu lafiya ƙwarai ; suna
miƙewa har shingen \VUgras{gonar shaƙatawa ta gari}
\textsuperscript{(17)}. \VUgras{Babbar hanya} \textsuperscript{(19)} wadda take
kai wa wannan gona tana wucewa a bayan \VUgras{bankin tanadi}
\textsuperscript{(18)} tana kuma kai wa filin \VUgras{babban coci}
\textsuperscript{(20)}.

%%K t11-c1-07-1 p
7. --- Dukan waɗannan gine-gine suna miƙewa a kan ƙaramin tudu kamar
matakai, inda babbar \VUgras{makarantar sakandarenmu}
\textsuperscript{(21)} take ma.

%%K t11-c1-07-2 p
Idan mutum ya sauka ta wata hanya mai ɗan gangara, wadda take haɗuwa
da Babbar Hanya, sai ya iso \VUgras{Kotu} \textsuperscript{(22)} wadda
\VUgras{lambun jama’a} \textsuperscript{(26)} mai daɗi yake miƙewa a
gabanta. Wannan \VUgras{filin lambu} ba babba ba ne ƙwarai, amma yana
yin marmaro na kore da sanyi a tsakiyar gari. Shi ya sa mutanen gari
suke zuwa gare shi da yawa, waɗanda suke son zuwa su zauna a ƙarƙashin
tantin \VUgras{gidan kofi} \textsuperscript{(25)}, domin su saurari makaɗan
sojoji waɗanda ranar Alhamis da ranar Lahadi suke kaɗawa a
\VUgras{rumfa} \textsuperscript{(27)} da aka ajiye tsakanin filayen ciyawa
da tsire-tsire.

%%K t11-c1-08-1 p
8. --- A kan ɗaya daga cikin waɗannan filayen ciyawa akwai
\VUgras{mutum-mutumi} \textsuperscript{(28)} na tagulla a tsaye, abin
tunawa da aka kafa domin tunawa da ɗaya daga cikin ƴan garinmu, wanda
ya yi wa garin haihuwarsa manyan hidimomi.

%%K t11-tit-3 sub
\VUsaut{4.0mm}
\VUpk{10.2pt}{40}{Kayan sufuri.}
\VUsaut{3.4mm}

%%K t11-c1-09-1 p
9. --- Har yanzu ba a haskaka garinmu da hasken lantarki ba, yana da
\VUgras{fitilun gas} \textsuperscript{(29)} kaɗai. Amma \VUgras{jaridun nan}
suna gwagwarmaya domin a inganta wannan tsarin haskakawa,
\VUgras{kamar} yadda aka riga aka inganta \VUgras{tsarin jan tiram.}
An maye gurbin \VUgras{tsofaffin} \VUgras{kekunan doki} da
\VUgras{tiram na lantarki} \textsuperscript{(31)} a aikace, waɗanda suke
jigilar matafiya da sauri daga ɗaya \VUgras{ƙarshen gari zuwa ɗaya,} a
kan farashi \VUgras{mai arha} ƙwarai.

%%K t11-c1-10-1 p
10. --- Kusa da Gidan Shari’a akwai \VUgras{Banki} \textsuperscript{(32)}
inda ake sayen takardun kasuwanci. \VUgras{Wakilan banki}
\textsuperscript{(33)} suna zuwa su karɓi bashi a gidajen mutane.

%%K t11-tit-4 sub
\VUsaut{4.0mm}
\VUpk{10.2pt}{40}{Fasaha da Ilimi.}
\VUsaut{3.4mm}

%%K t11-c1-11-1 p
11. --- Kyakkyawan ginin da yake tsaye a kusa shi ne \VUgras{gidan
kayan tarihi.} Yana ɗauke tare da \VUgras{ɗakin karatu}
\textsuperscript{(34)} na garin, da kyawawan \VUgras{tarin ilimin halitta}
\textsuperscript{(35)} da kuma \VUgras{gidan tarin zanuka}
\textsuperscript{(36)} mai kwarjini wanda yake yin \VUgras{wurin baje
koli} na dindindin mai kyau.

%%K t11-c1-11-2 p
Ana buɗe shi ga jama’a sau biyu a kowane mako, amma baƙi suna iya
ziyartarsa kowace rana idan sun ba \VUgras{mai gadi}
\textsuperscript{(37)} ɗan kyauta. \VUgras{Masu zane}
\textsuperscript{(38)} suna zuwa can domin su yi aiki. Muna ganin su
\VUgras{gaban} kayan zanensu, da allon fenti a hannu, suna kwafin
shahararrun zanuka.

%%K t11-c1-12-1 p
12. --- A kan tudu, a bayan gidan kayan tarihi, mun fara ganin
\VUgras{rumfar rufi} \textsuperscript{(40)} ta \VUgras{asibiti}
\textsuperscript{(39)} da gine-ginen \VUgras{makarantar horar da malamai}
\textsuperscript{(41)} inda aka horar da malaman makarantunmu na gari.

A kan filin Gidan Wasan Kwaikwayo akwai \VUgras{gidan waya}
\textsuperscript{(43)} da aka kafa a cikin \VUgras{gidan mutum}
\textsuperscript{(42)}.

%%K t11-tit-5 sub
\VUsaut{5.0mm}
\VUpk{11.4pt}{40}{Gobara.}
\VUsaut{3.0mm}

%%K t11-tit-6 sub
\VUpk{10.2pt}{40}{Kiran Gobara.}
\VUsaut{3.4mm}

%%K t11-c2-01-1 p
1. --- Labari mai ban tsoro ya bazu cikin garin : gobara ta tashi a
gidan wasan kwaikwayo yanzu ! A daidai lokacin da na iso
\VUgras{dandali} \textsuperscript{(96)}, tuni taron jama’a ya taru a kan
\VUgras{gefen hanya} \textsuperscript{(46)} da a kan \VUgras{babbar hanya}
\textsuperscript{(47)}. \VUgras{Ma’aikatan gidan waya}
\textsuperscript{(44)} da \VUgras{masu kai wasiƙa} \textsuperscript{(45)} suna
fitowa daga gidan wayar. \VUgras{Matuƙin tiram} \textsuperscript{(48)} ya
tilasta tsayar da abin hawansa domin ya bar \VUgras{motar asibiti}
\textsuperscript{(50)} ta gari ta wuce, wadda take zuwa da
\VUgras{likitan tiyata} \textsuperscript{(52)} da \VUgras{ma’aikacin jinya}
\textsuperscript{(51)} domin su kula da \VUgras{wanda ya ji rauni}
\textsuperscript{(53)} da aka kawo a kan \VUgras{shimfiɗar ɗauka}
\textsuperscript{(54)} kusa da \VUgras{rumfar jarida}
\textsuperscript{(55)}.

%%K t11-c2-02-1 p
2. --- Wata rundunar sojojin ƙasa, wadda take dawowa daga atisaye, ta
tsaya domin ta tabbatar da tsari tare da \VUgras{ƴan sandan doki na
gari} \textsuperscript{(61)}. \VUgras{Masu hawan dokin,} waɗanda dawakansu
suke tashi a tsaye, sun fi \VUgras{sojoji} \textsuperscript{(57)} ko
\VUgras{ƴan sanda} \textsuperscript{(63)} mayar da \VUgras{masu kallo}
\textsuperscript{(56)} baya. Ban da haka ƴan sandan suna cikin aiki
ƙwarai. Ɗaya daga cikinsu yana nuna wa \VUgras{ƴan sandan gari}
\textsuperscript{(64)} inda ya kamata a kai wani \VUgras{wanda ya sami
hatsari} \textsuperscript{(65)}, mai kashe gobara mai jaruntaka wanda ya
faɗi daga tsani.

%%K t11-c2-03-1 p
3. --- \VUgras{Jami’i} \textsuperscript{(66)} yana ƙayyade wurin da ya
kamata a ajiye \VUgras{famfon tururi} \textsuperscript{(67)} da yake zuwa.
Yana jiran wannan babbar na’ura, sai ya sa aka kafa \VUgras{famfon
hannu} \textsuperscript{(68)} da gaggawa, wanda doguwar
\VUgras{bututunsa} \textsuperscript{(69)} yake jefa ruwa kamar gudun ruwa a
kan wutar.

Masu kashe gobara shida suna motsa shi, alhali abokinsu, wanda yake
riƙe da \VUgras{bakin famfo} \textsuperscript{(70)}, yake jagorantar
\VUgras{yayyafin ruwa} \textsuperscript{(71)} zuwa tsakiyar gobarar.

%%K t11-c2-04-1 p
4. --- A ɗaya gefen gidan wasan kwaikwayo, an jingina babban
\VUgras{tsani} \textsuperscript{(72)}. Yana ba da damar masu kashe gobara su
kusanci harsunan wuta da suke fitowa cikin guguwar baƙin
\VUgras{hayaƙi} \textsuperscript{(75)}.

%%K t11-tit-7 sub
\VUsaut{5.0mm}
\VUpk{10.2pt}{40}{Ayyukan Jaruntaka.}
\VUsaut{3.4mm}

%%K t11-c2-05-1 p
5. --- \VUgras{Gobara} \textsuperscript{(74)} ta tashi a ɗakin hutu na
\VUgras{gidan wasan kwaikwayo} \textsuperscript{(73)}, alhali ana yin
maimaicin \VUgras{wasan waƙa} wanda \VUgras{wasansa} na farko zai
kasance gobe. Da sa’a \VUgras{masu kallo} \textsuperscript{(76)} kaɗan
kaɗai suke cikin zauren. Waɗannan masu tausayi sun gudu zuwa
\VUgras{baranda} \textsuperscript{(77)} inda \VUgras{masu kashe gobara}
\textsuperscript{(86)} masu jaruntaka suke zuwa su ceto su da tsani da
manyan igiyoyi.

%%K t11-c2-06-1 p
6. --- A ƙarƙashin \VUgras{shirayi} \textsuperscript{(78)}, inda
\VUgras{talla} \textsuperscript{(79)} yake sanar da wasan, muna ganin
\VUgras{makaɗa} \textsuperscript{(81)} na \VUgras{ƙungiyar makaɗa}
\textsuperscript{(80)} suna gudu, suna ɗauke da kayan kiɗansu :
\VUgras{bayolin} \textsuperscript{(81)}, da \VUgras{sarewa}
\textsuperscript{(82)}, da \VUgras{babban bayolin} \textsuperscript{(83)}, da
\VUgras{babbar busa} \textsuperscript{(84)}, da \VUgras{ganga}
\textsuperscript{(85)}, da sauransu. Ana ƙoƙarin ceto ɓangaren
\VUgras{kayan gida} \textsuperscript{(87)}.

%%K t11-c2-07-1 p
7. --- \VUgras{Ƴan wasa maza} \textsuperscript{(89)} da \VUgras{ƴan wasa
mata} \textsuperscript{(88)}, da \VUgras{mawaƙa maza} da \VUgras{mawaƙa
mata} sun tilasta gudu da \VUgras{kayan wasansu}
\textsuperscript{(90)}. Mawaƙi mai murya yana bayyana wa \VUgras{ɗan
jarida} \textsuperscript{(92)} yadda abin ya faru. \VUgras{Mai bayar da
rahoto} ya gaggauta zuwa tun daga farko tare da masu mulki :
\VUgras{gwamna} \textsuperscript{(91)}, da \VUgras{magajin gari}
\textsuperscript{(93)}, da \VUgras{kwamishinan ƴan sanda}
\textsuperscript{(94)} da \VUgras{jami’in shari’a} \textsuperscript{(95)},
waɗanda za su yi bincike domin su yanke hukunci a kan alhakin.

\VUsaut{6.0mm}
\VUfilet{22mm}
